\section{Glossary of terms}\label{app:glossary}

\begin{description}[nosep,leftmargin=0em,style=unboxed,font=\normalfont\itshape]
\item[AIS.] The Automatic Identification System: a VHF broadcast every vessel over
  300 gross tonnes must transmit, carrying its identity (MMSI), position, speed and
  course, and less often its draught, destination and ETA. Terrestrial receivers
  hear it within roughly 40 nautical miles of shore; satellites hear it at sea.
\item[MMSI.] Maritime Mobile Service Identity, the nine-digit number that keys AIS
  messages to a vessel. It can change on reflagging, which is why the fleet list is
  keyed on the IMO hull number and resolved to an MMSI.
\item[FSRU.] Floating storage and regasification unit, a converted carrier moored
  permanently as an import terminal.
\item[Laden, ballast.] A carrier is laden when carrying cargo and in ballast when
  empty; Equation~\eqref{eq:laden} infers the state from draught.
\item[Leg, visit, queue.] The three event pairings: departure to next arrival
  (a voyage), mooring to departure (berth occupancy), first anchorage entry to
  mooring (the wait to berth).
\item[Physical, knowable.] The two bases every signal is built on: with hindsight,
  and using only what was observable on the day (Section~\ref{sec:panel}).
\item[Vintage.] The version of a dataset as it existed on a given date. A feature
  has a vintage hazard when the table it comes from did not exist at the as-of date
  even though its rows carry earlier timestamps.
\item[Persistence null, climatology null.] Forecasting next week's count as last
  week's count, and as the trailing four-week mean, respectively.
\item[MAE, skill, excess.] Mean absolute error; the fractional MAE improvement over
  a null, Equation~\eqref{eq:skill}; and the fractional MAE deterioration relative
  to a null, Equation~\eqref{eq:excess}.
\item[80\% coverage.] The share of realised outcomes that fell inside the model's
  central 80\% predictive interval. A calibrated model scores about 80\%.
\item[Walk-forward.] Refitting the model at every step on data strictly before that
  step, so that every forecast is out of sample. Contrast with cross-validation,
  which shuffles time.
\item[Purge and embargo.] Dropping training rows whose target window overlaps the
  test date, and one further row as a margin (Section~\ref{sec:partb}).
\item[HAC, Newey--West.] A heteroskedasticity- and autocorrelation-consistent
  standard error, Equation~\eqref{eq:nw}, which corrects for the serial correlation
  that overlapping targets induce.
\item[Bonferroni.] Dividing the significance level by the number of tests so that
  the family-wise false-positive rate stays at the nominal level.
\item[FWL, partial effect, partial \Rsq.] The Frisch--Waugh--Lovell construction of
  a coefficient net of controls, Equation~\eqref{eq:fwl}, and the share of the
  residual variance that the signal explains.
\item[Discovery, holdout.] The sample on which a hypothesis is graded, and the
  final stretch left untouched until it has passed.
\item[EWMA.] Exponentially weighted moving average, $\hat{x}_t = \alpha y_t + (1-\alpha)\hat{x}_{t-1}$.
\item[BOCPD.] Bayesian online change-point detection, Equation~\eqref{eq:bocpd}.
\item[CRPS.] Continuous ranked probability score, a proper score for a full
  predictive distribution that reduces to absolute error for a point forecast.
\item[Capture rate.] Observed cargoes divided by the EIA-implied count,
  Equation~\eqref{eq:capture}.
\item[Coverage seam.] The date after which only the thin live feed carries the
  signals (Section~\ref{sec:seam}).
\end{description}
